%% kapa: exam=1 round=27 year=2016 period=2 form=A subject=appraisal_law subject_ko=감정평가관계법규 seq=29 number=29
%% source: Q-Net 감정평가사 기출문제 (cst00309) artlSeq=5169402 file=extracted/1008/5169402/제27회 감정평가사 1차 문제/제27회 감정평가사 2교시문제(A형).pdf
%% conversion: kordoc 4.24.51 -> tools/gen1.py (auto); review=proofread-sonnet; flags=none
\begin{question}{29}
공간정보의 구축 및 관리 등에 관한 법령상 지적공부와 부동산등기부에 관한 설명으로 옳지 않은 것은?
\begin{choices}
  \item 지적소관청은 지적공부의 등록사항중 지적도 및 임야도에 등록된 필지가 면적의 증감 없이 경계의 위치만 잘못된 경우를 발견하면 직권으로 조사ㆍ측량하여 정정할 수 있다.
  \item 지적소관청 소속 공무원이 지적공부와 부동산등기부의 부합 여부를 확인하기 위하여 등기사항증명서의 발급을 신청하는 경우 그 수수료를 감경할 수 있다.
  \item 행정구역의 명칭이 변경되었으면 지적공부에 등록된 토지의 소재는 새로운 행정구역의 명칭으로 변경된 것으로 본다.
  \item 지적공부에 신규등록하는 토지의 소유자에 관한 사항은 지적소관청이 직접 조사하여 등록한다.
  \item 「국유재산법」상 중앙관서의 장이 소유자 없는 부동산에 대한 소유자 등록을 신청하는 때에 지적소관청은 지적공부에 해당 토지의 소유자가 등록되지 아니한 경우에만 등록할 수 있다.
\end{choices}
\end{question}
\answer{2}
